% =============================================================================
% Communication Strategy - Method Section (for Edge AI Chapter)
% =============================================================================

\subsection{Network Architecture}
\label{sec:network-architecture}

The Drillbot system requires reliable communication between the Raspberry Pi (hosting the RAG-based assistant) and the Control System PC (running the operator GUI). This section describes the network architecture and the rationale for design decisions.

\subsubsection{System Components}

The communication system consists of three networked components connected via an Ethernet switch:

\begin{itemize}
    \item \textbf{Raspberry Pi:} Hosts the Drillbot API server (FastAPI) and the local LLM (Ollama with Qwen 2.5 1.5B).
    \item \textbf{Control System PC:} Runs the operator GUI with the Drillbot chat interface.
    \item \textbf{Trajectory System:} Dedicated workstation for path planning (pre-existing infrastructure).
\end{itemize}

\subsubsection{Network Topology}

All rig components communicate over a dedicated Ethernet network using static IP addresses on the \texttt{10.10.10.0/24} subnet. Figure~\ref{fig:network-topology} illustrates the network topology.

\begin{figure}[htbp]
    \centering
    \begin{tikzpicture}[
        node distance=2cm,
        device/.style={rectangle, draw, minimum width=2.5cm, minimum height=1cm, align=center},
        switch/.style={rectangle, draw, minimum width=3cm, minimum height=0.8cm, fill=gray!20}
    ]
        % Devices
        \node[device] (pi) {Raspberry Pi\\{\small 10.10.10.20}};
        \node[device, right=of pi] (control) {Control PC\\{\small 10.10.10.5}};
        \node[device, right=of control] (traj) {Trajectory\\{\small 10.10.10.4}};

        % Switch
        \node[switch, below=1.5cm of control] (switch) {Ethernet Switch};

        % Connections
        \draw (pi) -- (switch);
        \draw (control) -- (switch);
        \draw (traj) -- (switch);

        % Network label
        \node[below=0.3cm of switch, font=\small\itshape] {Isolated rig network (10.10.10.0/24)};
    \end{tikzpicture}
    \caption{Network topology for the Drillbot communication system. All devices are connected via an unmanaged Ethernet switch on an isolated network segment.}
    \label{fig:network-topology}
\end{figure}

\subsubsection{IP Address Allocation}

Table~\ref{tab:ip-allocation} presents the static IP address allocation for the rig network.

\begin{table}[htbp]
    \centering
    \caption{Static IP address allocation on the rig Ethernet network.}
    \label{tab:ip-allocation}
    \begin{tabular}{llll}
        \toprule
        \textbf{Device} & \textbf{IP Address} & \textbf{Interface} & \textbf{Configuration} \\
        \midrule
        Raspberry Pi & 10.10.10.20 & eth0 & NetworkManager (static) \\
        Control PC & 10.10.10.5 & Ethernet & Windows static IP \\
        Trajectory System & 10.10.10.4 & Ethernet & Static \\
        \bottomrule
    \end{tabular}
\end{table}

\subsubsection{Communication Protocol}

The Drillbot system uses HTTP/REST for communication between the GUI and the Raspberry Pi. Table~\ref{tab:api-endpoints} summarizes the API endpoints.

\begin{table}[htbp]
    \centering
    \caption{Drillbot API endpoints hosted on the Raspberry Pi.}
    \label{tab:api-endpoints}
    \begin{tabular}{llll}
        \toprule
        \textbf{Endpoint} & \textbf{Method} & \textbf{Purpose} & \textbf{Port} \\
        \midrule
        \texttt{/health} & GET & Service health check & 8000 \\
        \texttt{/ask} & POST & Submit question to RAG system & 8000 \\
        \texttt{/ask/stream} & POST & Streaming response variant & 8000 \\
        \bottomrule
    \end{tabular}
\end{table}

The choice of HTTP over HTTPS was deliberate: the rig network is physically isolated with no internet connectivity, and only authorized equipment is connected. In this closed environment, the overhead of TLS encryption provides no practical security benefit while adding certificate management complexity.

\subsubsection{Rationale for Dedicated Ethernet}

The dedicated Ethernet network was chosen over the university WiFi infrastructure (eduroam) for the following reasons:

\begin{enumerate}
    \item \textbf{Static IP addressing:} The university wireless network uses DHCP with dynamic address allocation. IP addresses change upon reconnection or lease expiration, requiring constant reconfiguration of client systems. According to NTNU IT policy, static IP addresses cannot be assigned on wireless networks\footnote{NTNU IT: ``Merk at det kun kan tildeles fast IP-adresse i kablet nettverk, vi kan ikke sette fast IP-adresse i NTNUs tr{\aa}dl{\o}se nett.''}.

    \item \textbf{Reliability:} Wired Ethernet provides deterministic latency (<1\,ms) compared to WiFi (2--10\,ms typical). For a control system application, predictable communication timing is essential.

    \item \textbf{Isolation:} The rig network operates independently of university infrastructure, eliminating external dependencies and potential interference from other network traffic.

    \item \textbf{Compatibility:} The Control PC and Trajectory System already used the \texttt{10.10.10.0/24} subnet; integrating the Raspberry Pi into this existing network required minimal changes to the established infrastructure.
\end{enumerate}

\subsubsection{Service Configuration}

The Drillbot API runs as a systemd service on the Raspberry Pi, ensuring automatic startup on boot and automatic restart on failure. Listing~\ref{lst:systemd-service} shows the service configuration.

\begin{lstlisting}[
    caption={Systemd service configuration for the Drillbot API.},
    label={lst:systemd-service},
    language=bash,
    basicstyle=\ttfamily\small,
    frame=single
]
[Unit]
Description=Drillbot RAG API
After=network.target

[Service]
Type=simple
User=drillbotics
WorkingDirectory=/home/drillbotics/drillbot
ExecStart=/home/drillbotics/drillbot/venv/bin/uvicorn \
    app.api:app --host 0.0.0.0 --port 8000
Restart=always
RestartSec=5

[Install]
WantedBy=multi-user.target
\end{lstlisting}

This configuration enables headless operation of the Raspberry Pi without requiring a monitor, keyboard, or manual intervention after power-on.


% =============================================================================
% Appendix Section - WiFi Configuration (Non-Ideal Solution)
% =============================================================================

% Add this to your appendix chapter

\section{WiFi Communication Attempt}
\label{appendix:wifi-communication}

This appendix documents the initial attempt to use the university WiFi network (eduroam) for Drillbot communication. This approach was ultimately abandoned in favor of the dedicated Ethernet network described in Section~\ref{sec:network-architecture}.

\subsection{Initial Configuration}

The Raspberry Pi was initially configured to connect to the university WiFi network:

\begin{itemize}
    \item \textbf{Network:} eduroam
    \item \textbf{Protocol:} WPA2-Enterprise (EAP-PEAP)
    \item \textbf{IP Assignment:} DHCP (dynamic)
\end{itemize}

The GUI configuration pointed to the Pi's current WiFi IP address:

\begin{lstlisting}[language=Python, basicstyle=\ttfamily\small, frame=single]
# config.py - WiFi configuration
PI_HOST = "10.22.60.77"  # Dynamic IP - changes on reconnect
PI_PORT = 8000
\end{lstlisting}

\subsection{Problems Encountered}

\subsubsection{Dynamic IP Address Assignment}

The primary issue was the instability of the IP address. Table~\ref{tab:ip-changes} shows the IP addresses observed across multiple connection sessions.

\begin{table}[htbp]
    \centering
    \caption{Observed IP address changes on university WiFi.}
    \label{tab:ip-changes}
    \begin{tabular}{ll}
        \toprule
        \textbf{Session} & \textbf{Assigned IP Address} \\
        \midrule
        First connection & 10.22.61.40 \\
        Second connection & 10.22.61.83 \\
        Third connection & 10.22.60.77 \\
        Fourth connection & 10.22.60.223 \\
        \bottomrule
    \end{tabular}
\end{table}

Each IP change required:
\begin{enumerate}
    \item Determining the new IP address on the Raspberry Pi (\texttt{hostname -I})
    \item Updating configuration files on the Control PC
    \item Updating configuration files in the GUI codebase
    \item Restarting affected services
\end{enumerate}

\subsubsection{University Policy Constraints}

NTNU IT policy explicitly prohibits static IP assignment on wireless networks:

\begin{quote}
    ``Adresser deles ut dynamisk til maskiner som er koblet til nettverket. Hvis en maskin kobles fra nettet over lang tid kan adressen som var tildelt bli delt ut til andre. Derfor vil det variere hvilken adresse en maskin har over tid, selv om du ikke flytter p{\aa} maskinen.'' \\
    \textit{(Addresses are distributed dynamically to machines connected to the network. If a machine is disconnected for a long time, the assigned address may be given to others. Therefore, the address a machine has will vary over time, even if the machine is not moved.)}
\end{quote}

\subsubsection{Operational Impact}

The dynamic IP issue manifested as:
\begin{itemize}
    \item \textbf{Connection failures:} The GUI could not reach the Drillbot API after IP changes
    \item \textbf{Manual intervention required:} Each reconnection required operator action
    \item \textbf{Unpredictable timing:} IP changes occurred without warning
\end{itemize}

\subsection{Workarounds Attempted}

Several workarounds were evaluated:

\begin{enumerate}
    \item \textbf{mDNS/Avahi:} Attempted to use \texttt{raspberrypi.local} hostname resolution. Failed due to enterprise network restrictions blocking multicast DNS.

    \item \textbf{Reserved DHCP lease:} Requested a static DHCP reservation from NTNU IT. Denied per university policy for wireless devices.

    \item \textbf{VPN tunnel:} Considered establishing a VPN connection. Rejected due to added complexity and introduction of additional failure points.
\end{enumerate}

\subsection{Conclusion}

The WiFi-based approach proved unsuitable for a production control system due to:
\begin{itemize}
    \item Unpredictable IP address changes
    \item University policy prohibiting static wireless IP assignment
    \item Requirement for manual reconfiguration after each network event
\end{itemize}

The dedicated Ethernet network (Section~\ref{sec:network-architecture}) resolved all of these issues by providing static, predictable addressing on an isolated network segment under full project control.
